%**************************************************************************************
% License:
% CC BY-NC-SA 4.0 (http://creativecommons.org/licenses/by-nc-sa/4.0/)
%**************************************************************************************

\documentclass[notes]{beamer}

\mode<presentation> {

\usetheme{Madrid}

% Burnt orange
\definecolor{burntorange}{rgb}{0.8, 0.33, 0.0}
\colorlet{beamer@blendedblue}{burntorange}
% Pale yellow
\definecolor{paleyellow}{rgb}{1.0, 1.0, 0.953}
\setbeamercolor{background canvas}{bg=white}
% Secondary and tertiary palette
\setbeamercolor*{palette secondary}{use=structure,fg=white,bg=burntorange!80!black}
\setbeamercolor*{palette tertiary}{use=structure,fg=white,bg=burntorange!60!black}

% To remove the navigation symbols from the bottom of all slides uncomment this line
%\setbeamertemplate{navigation symbols}{}
}

\usepackage{amsmath}
\DeclareMathOperator*{\argmin}{arg\,min}
\DeclareMathOperator*{\argmax}{arg\,max}
\usepackage{bm}
\usepackage{booktabs} % Allows the use of \toprule, \midrule and \bottomrule in tables
\usepackage{breqn}
\usepackage{cleveref}
\usepackage{graphicx} % for figures
\usepackage[labelsep=space,tableposition=top]{caption}
\renewcommand{\figurename}{Fig.}
\usepackage{caption,subcaption}% http://ctan.org/pkg/{caption,subcaption}
\usepackage{xcolor}

%----------------------------------------------------------------------------------------
%	TITLE PAGE
%----------------------------------------------------------------------------------------
\title[Exploratory Data Analysis]{Exploratory Data Analysis for Liquefaction Assessment}
\author{Krishna Kumar}
\institute[UT Austin]
{
University of Texas at Austin \\
\medskip
\textit{\url{krishnak@utexas.edu}}
}
\date{}

\begin{document}

\begin{frame}
\titlepage
\end{frame}

\begin{frame}
 \frametitle{Overview}
 \tableofcontents
\end{frame}

%----------------------------------------------------------------------------------------
% slides
%----------------------------------------------------------------------------------------

\section{Why EDA Matters}

%------------------------------------------------
\begin{frame}
\frametitle{The Foundation of Machine Learning}

\textbf{Before building models, understand your data!}

\mode<beamer>{
\textbf{EDA answers:}
\begin{itemize}
\item What does the data tell us about the physics?
\item Are there quality issues (missing, errors, outliers)?
\item Which features predict lateral spreading?
\item How should we preprocess?
\end{itemize}
}
\mode<handout>{
\vspace{3cm}
}

\vspace{1em}

\textbf{Gilbert Strang's approach:} \\
\textit{"See the structure first, then calculate."}

\end{frame}

%------------------------------------------------
\begin{frame}
\frametitle{Liquefaction and Lateral Spreading}

\textbf{Physical process:}
\begin{enumerate}
\item Earthquake shaking
\item Pore pressure increases
\item Effective stress drops: $\sigma' = \sigma - u$
\item Soil loses strength
\item Lateral displacement occurs
\end{enumerate}

\vspace{1em}

\textbf{Dataset:} Durante \& Rathje (2021)
\begin{itemize}
\item 7,291 case histories
\item Target: >0.3m displacement = spreading
\end{itemize}

\end{frame}

%------------------------------------------------
\section{Understanding Distributions}

%------------------------------------------------
\begin{frame}
\frametitle{The Key Question}

\textbf{Can we distinguish spreading from non-spreading sites?}

\mode<beamer>{
\textbf{Approach:} Compare distributions by class

\begin{itemize}
\item \textbf{Overlapping distributions:} Feature not useful
\item \textbf{Separated distributions:} Feature is predictive
\end{itemize}

\vspace{1em}

\textbf{Best visualization:} Overlaid histograms
\begin{itemize}
\item See both classes simultaneously
\item Directly compare shapes and locations
\item Immediate intuition about separability
\end{itemize}
}
\mode<handout>{
\vspace{3cm}
}

\textbf{This is more informative than looking at all data together!}

\end{frame}

%------------------------------------------------
\begin{frame}
\frametitle{Reading Overlaid Histograms}

\textbf{What to look for:}

\mode<beamer>{
\begin{enumerate}
\item \textbf{Separation:} How far apart are the peaks?
  \begin{itemize}
  \item Far apart $\to$ strong predictor
  \item Overlapping $\to$ weak predictor
  \end{itemize}

\item \textbf{Direction:} Which class has higher values?
  \begin{itemize}
  \item GWD: No-spread sites have higher values (deeper water table = safer)
  \item PGA: Spread sites have higher values (more shaking = more risk)
  \end{itemize}

\item \textbf{Overlap area:} How much do they overlap?
  \begin{itemize}
  \item Small overlap $\to$ easy to classify
  \item Large overlap $\to$ difficult boundary
  \end{itemize}
\end{enumerate}
}
\mode<handout>{
\vspace{3cm}
}

\end{frame}

%------------------------------------------------
\section{Quartiles: Robust Statistics}

%------------------------------------------------
\begin{frame}
\frametitle{Why Quartiles Matter}

\textbf{Problem:} Mean and std are sensitive to outliers

\mode<beamer>{
\textbf{Example:} Water depths (m): \{1.5, 1.8, 2.0, 2.1, 2.3, 15.0\}
\begin{itemize}
\item Mean = 4.12 m (misleading!)
\item Median = 2.05 m (representative!)
\end{itemize}

\vspace{1em}

\textbf{Quartiles divide data into four equal parts:}
\begin{itemize}
\item Q1 (25th percentile): 25\% below
\item Q2 (Median, 50th percentile): middle value
\item Q3 (75th percentile): 75\% below
\item IQR = Q3 - Q1 (middle 50\%)
\end{itemize}
}
\mode<handout>{
\vspace{3cm}
}

\textbf{Box plots visualize the five-number summary}

\end{frame}

%------------------------------------------------
\begin{frame}
\frametitle{Quartiles for Engineering Decisions}

\textbf{Example: Ground Water Depth (GWD)}

\mode<beamer>{
From our data:
\begin{itemize}
\item Q1 = 1.64 m, Median = 1.98 m, Q3 = 2.46 m
\item IQR = 0.82 m
\end{itemize}

\vspace{0.5em}

\textbf{Engineering interpretation:}
\begin{itemize}
\item 50\% of sites: water table at 1.64-2.46 m depth
\item Sites with GWD < Q1 (1.64 m): \textbf{High risk} zone
\item Sites with GWD > Q3 (2.46 m): \textbf{Lower risk} zone
\end{itemize}

\vspace{0.5em}

\textbf{Action:} Quartiles create natural risk categories!
}
\mode<handout>{
\vspace{3cm}
}

\end{frame}

%------------------------------------------------
\section{Outlier Detection}

%------------------------------------------------
\begin{frame}
\frametitle{The Outlier Dilemma}

\textbf{In geotechnical engineering, outliers can be:}

\mode<beamer>{
\begin{enumerate}
\item \textbf{Measurement errors:} Remove
  \begin{itemize}
  \item Negative water table depth
  \item Physically impossible values
  \end{itemize}

\item \textbf{Extreme events:} Keep!
  \begin{itemize}
  \item Major earthquakes
  \item Unusual but valid site conditions
  \end{itemize}
\end{enumerate}

\vspace{1em}

\textbf{Never blindly remove outliers!}
\begin{itemize}
\item They might be your most important data
\item Use engineering judgment + statistics
\end{itemize}
}
\mode<handout>{
\vspace{3cm}
}

\end{frame}

%------------------------------------------------
\begin{frame}
\frametitle{IQR Method (Most Robust)}

\textbf{Rule:} Values beyond these bounds are outliers

\mode<beamer>{
\begin{align*}
\text{Lower bound} &= Q1 - 1.5 \times \text{IQR} \\
\text{Upper bound} &= Q3 + 1.5 \times \text{IQR}
\end{align*}

\textbf{Why 1.5×IQR?}
\begin{itemize}
\item John Tukey's empirical rule
\item For normal distribution: captures $\sim$99.3\% of data
\item Balances sensitivity vs over-flagging
\end{itemize}

\vspace{0.5em}

\textbf{Advantages:}
\begin{itemize}
\item Works for any distribution shape
\item Not fooled by outliers themselves
\item Easy to explain to engineers
\end{itemize}
}
\mode<handout>{
\vspace{3cm}
}

\end{frame}

%------------------------------------------------
\section{Z-Scores: Two Uses}

%------------------------------------------------
\begin{frame}
\frametitle{What is a Z-Score?}

\textbf{Z-score = how many standard deviations from the mean}

\mode<beamer>{
\begin{equation*}
z = \frac{x - \mu}{\sigma}
\end{equation*}

\textbf{Interpretation:}
\begin{itemize}
\item $z = 0$: At the mean
\item $z = 1$: One std above mean
\item $z = -2$: Two std below mean
\item $|z| > 3$: Outlier (only 0.3\% of normal data)
\end{itemize}

\vspace{1em}

\textbf{68-95-99.7 Rule:}
\begin{itemize}
\item 68\% within $\pm 1\sigma$
\item 95\% within $\pm 2\sigma$
\item 99.7\% within $\pm 3\sigma$
\end{itemize}
}
\mode<handout>{
\vspace{3cm}
}

\end{frame}

%------------------------------------------------
\begin{frame}
\frametitle{Use 1: Outlier Detection}

\textbf{Example: Ground Water Depth}

\mode<beamer>{
Data: Mean $\mu = 2.08$ m, Std $\sigma = 0.64$ m

\vspace{0.5em}

\begin{tabular}{lcc}
\toprule
\textbf{Value} & \textbf{Z-score} & \textbf{Assessment} \\
\midrule
2.08 m & 0 & Average \\
2.72 m & 1 & Normal \\
3.36 m & 2 & Unusual \\
6.00 m & 6.13 & \textbf{Outlier!} \\
\bottomrule
\end{tabular}

\vspace{0.5em}

\textbf{Threshold:} $|z| > 3$ for outliers
}
\mode<handout>{
\vspace{3cm}
}

\end{frame}

%------------------------------------------------
\begin{frame}
\frametitle{Use 2: Standardization (More Important!)}

\textbf{Problem:} Features have different scales
\begin{itemize}
\item GWD: 0.4 - 6.0 m
\item PGA: 0.3 - 0.6 g
\end{itemize}

\vspace{1em}

\textbf{Solution:} Standardize using z-scores

\mode<beamer>{
Transform each feature:
\begin{equation*}
x_{\text{standardized}} = \frac{x - \mu}{\sigma}
\end{equation*}

\textbf{Result:} All features have mean = 0, std = 1

\vspace{1em}

\textbf{Why this matters for ML:}
\begin{itemize}
\item Neural networks train better
\item Distance-based methods (KNN, SVM) need equal scales
\item Prevents large-value features from dominating
\end{itemize}
}
\mode<handout>{
\vspace{3cm}
}

\end{frame}

%------------------------------------------------
\begin{frame}
\frametitle{IQR vs Z-Score: When to Use Each}

\begin{center}
\begin{tabular}{lll}
\toprule
\textbf{Method} & \textbf{Use For} & \textbf{Advantage} \\
\midrule
IQR & Outlier detection & Robust \\
Z-score & Outlier detection & Assumes normality \\
Z-score & \textbf{Standardization} & \textbf{ML preprocessing} \\
\bottomrule
\end{tabular}
\end{center}

\vspace{1em}

\textbf{Recommendation:}
\begin{itemize}
\item Detect outliers with IQR
\item Standardize features with z-scores for neural networks
\item Tree-based models don't need standardization
\end{itemize}

\end{frame}

%------------------------------------------------
\section{Correlation and Significance}

%------------------------------------------------
\begin{frame}
\frametitle{Correlation Coefficient}

\textbf{Measures linear relationship strength}

\mode<beamer>{
\textbf{Properties:}
\begin{itemize}
\item Range: $r \in [-1, 1]$
\item $r = 1$: Perfect positive correlation
\item $r = -1$: Perfect negative correlation
\item $r = 0$: No linear correlation
\end{itemize}

\vspace{0.5em}

\textbf{Interpretation for liquefaction:}
\begin{itemize}
\item GWD vs Target: Negative (shallow water = more risk)
\item PGA vs Target: Positive (more shaking = more risk)
\item L vs Target: Negative (closer to river = more risk)
\end{itemize}
}
\mode<handout>{
\vspace{3cm}
}

\textbf{Warning:} Correlation $\neq$ Causation!

\end{frame}

%------------------------------------------------
\begin{frame}
\frametitle{Statistical Significance Testing}

\textbf{Question:} Are distributions truly different, or just random chance?

\mode<beamer>{
\textbf{Mann-Whitney U Test:}
\begin{itemize}
\item Compares two distributions
\item Non-parametric (no normality assumption)
\item Returns p-value
\end{itemize}

\vspace{1em}

\textbf{P-value interpretation:}
\begin{itemize}
\item $p < 0.001$: Extremely significant (***)
\item $p < 0.01$: Very significant (**)
\item $p < 0.05$: Significant (*)
\item $p \geq 0.05$: Not significant
\end{itemize}

\vspace{1em}

\textbf{For our data:} GWD, PGA, and L are highly significant (p < 0.001)
}
\mode<handout>{
\vspace{3cm}
}

\end{frame}

%------------------------------------------------
\section{Data Quality}

%------------------------------------------------
\begin{frame}
\frametitle{Quick Quality Checks}

\textbf{Before modeling, verify:}

\mode<beamer>{
\begin{enumerate}
\item \textbf{Missing values:}
  \begin{itemize}
  \item How many? Which features?
  \item Delete (<5\%) or impute (>5\%)?
  \end{itemize}

\item \textbf{Physical plausibility:}
  \begin{itemize}
  \item Negative values where impossible
  \item Values outside engineering bounds
  \end{itemize}

\item \textbf{Duplicates:}
  \begin{itemize}
  \item Exact copies (data entry errors)?
  \item Same site, multiple events (valid)?
  \end{itemize}

\item \textbf{Class balance:}
  \begin{itemize}
  \item 50-50 split? 90-10?
  \item Severe imbalance needs special handling
  \end{itemize}
\end{enumerate}
}
\mode<handout>{
\vspace{3cm}
}

\end{frame}

%------------------------------------------------
\section{Key Takeaways}

%------------------------------------------------
\begin{frame}
\frametitle{Essential Concepts}

\textbf{Must-know statistical tools:}

\mode<beamer>{
\begin{enumerate}
\item \textbf{Quartiles (Q1, Median, Q3):}
  \begin{itemize}
  \item Robust alternative to mean/std
  \item IQR = middle 50\% of data
  \end{itemize}

\item \textbf{Z-scores:} $z = (x - \mu) / \sigma$
  \begin{itemize}
  \item Outlier detection: $|z| > 3$
  \item Standardization for ML
  \end{itemize}

\item \textbf{Correlation:}
  \begin{itemize}
  \item Strength of linear relationship
  \item Not causation!
  \end{itemize}

\item \textbf{Statistical tests:}
  \begin{itemize}
  \item Confirm differences aren't random
  \item P-value < 0.05 = significant
  \end{itemize}
\end{enumerate}
}
\mode<handout>{
\vspace{3cm}
}

\end{frame}

%------------------------------------------------
\begin{frame}
\frametitle{Insights from Liquefaction Data}

\textbf{What we learned:}

\mode<beamer>{
\begin{enumerate}
\item \textbf{Most important features:}
  \begin{itemize}
  \item GWD: Strong negative correlation
  \item PGA: Strong positive correlation
  \item L: Moderate negative correlation
  \end{itemize}

\item \textbf{Data quality:} Excellent
  \begin{itemize}
  \item No missing values
  \item All values physically plausible
  \item Minimal duplicates
  \end{itemize}

\item \textbf{Class balance:} ~60/40 split
  \begin{itemize}
  \item Reasonable balance
  \item Use stratified sampling
  \end{itemize}
\end{enumerate}

\vspace{1em}

\textbf{Ready for machine learning!}
}
\mode<handout>{
\vspace{3cm}
}

\end{frame}

%------------------------------------------------
\begin{frame}
\frametitle{ML Modeling Strategy}

\textbf{Based on EDA:}

\mode<beamer>{
\begin{enumerate}
\item \textbf{Preprocessing:}
  \begin{itemize}
  \item Standardize (z-scores) for neural networks
  \item Raw features for tree-based models
  \item Stratified train-test split
  \end{itemize}

\item \textbf{Model selection:}
  \begin{itemize}
  \item Start with Random Forest (robust, interpretable)
  \item Try XGBoost (often best performance)
  \item Neural networks with standardized inputs
  \end{itemize}

\item \textbf{Evaluation metrics:}
  \begin{itemize}
  \item Multiple metrics: Accuracy, Precision, Recall, F1, AUC
  \item Confusion matrix
  \item Consider cost of errors in engineering context
  \end{itemize}
\end{enumerate}
}
\mode<handout>{
\vspace{3cm}
}

\end{frame}

%------------------------------------------------
\begin{frame}
\frametitle{Words of Wisdom}

\begin{center}
\large

\textit{"All models are wrong, but some are useful."} \\
\vspace{0.3em}
--- George Box

\vspace{2em}

\textit{"See the structure first, then calculate."} \\
\vspace{0.3em}
--- Gilbert Strang

\vspace{2em}

\textit{"Connect to the application - numbers must have physical meaning."}
\end{center}

\vspace{1em}

\textbf{The goal is insight, not just numbers!}

\end{frame}

%------------------------------------------------
\begin{frame}
\frametitle{Common Mistakes to Avoid}

\mode<beamer>{
\begin{enumerate}
\item \textbf{Blindly trusting summary statistics}
  \begin{itemize}
  \item Always visualize!
  \end{itemize}

\item \textbf{Removing outliers without investigation}
  \begin{itemize}
  \item Might be your most important data
  \end{itemize}

\item \textbf{Confusing correlation with causation}
  \begin{itemize}
  \item Ice cream sales correlate with drowning... both caused by summer!
  \end{itemize}

\item \textbf{Ignoring class imbalance}
  \begin{itemize}
  \item 99\% accuracy meaningless for 1\% rare class
  \end{itemize}

\item \textbf{Skipping domain knowledge check}
  \begin{itemize}
  \item Does this make engineering sense?
  \end{itemize}
\end{enumerate}
}
\mode<handout>{
\vspace{3cm}
}

\end{frame}

%------------------------------------------------
\begin{frame}
\frametitle{Summary}

\textbf{EDA is the foundation of successful ML}

\vspace{1em}

\mode<beamer>{
\textbf{We accomplished:}
\begin{itemize}
\item Understood data structure and quality
\item Identified predictive features
\item Compared distributions by class
\item Detected and evaluated outliers
\item Gained physical insight into liquefaction
\item Made informed preprocessing decisions
\end{itemize}

\vspace{1em}

\textit{"No amount of sophisticated modeling can compensate for poor understanding of your data."}
}
\mode<handout>{
\vspace{4cm}
}

\end{frame}

%------------------------------------------------
\begin{frame}
\frametitle{Questions?}

\centering
\Large Thank you!

\vspace{2cm}

\textbf{Contact:} \\
Krishna Kumar \\
\textit{krishnak@utexas.edu} \\
University of Texas at Austin

\vspace{1cm}

\textbf{Remember:} \\
\textit{"Understanding the data is understanding the physics."}

\end{frame}

\end{document}
